% GENERATED FILE. Edit canonical lessons, then run npm run generate:books.
% canonical-source-hash: fnv1a64:b193e9b7513becc0
% canonical-lessons: RU-C312-pozdravlyat, RU-C312-pozdravleniye, RU-C312-pozhelaniye, RU-C312-zhelat, RU-C312-tost

\chapter{Congratulations and wishes}
\label{ch:ru-congratulations-and-wishes}
\hlchaptermodality{312}

\begin{chapteropening}
\textbf{By the end of this chapter:} \emph{I can congratulate someone, talk about greetings and wishes, wish someone well and raise a toast.}

Complete the last lesson of chapter 312: I can congratulate someone, talk about greetings and wishes, wish someone well and raise a toast.
\end{chapteropening}

\section[\texorpdfstring{\ru{поздравлять} (pozdravlyát)}{поздравлять (pozdravlyát)}]{\ru{поздравлять} (pozdravlyát) — to congratulate}

\label{lesson:RU-C312-pozdravlyat}



\begin{quote}
Before the new one: say the Russian for solemn, then the Russian for festive.
\end{quote}

\subsection*{You'll want to know: \ru{поздравлять}}
\textbf{\ru{поздравлять}} (\emph{pozdravlyát}) — \textquotedblleft{}to congratulate\textquotedblright{}.

\begin{quote}
\ru{Друзья} \ru{поздравили} \ru{меня} \ru{с} \ru{днём} \ru{рождения}.
\end{quote}

\begin{quote}
\emph{Druzyá pozdrávili menyá s dnyom rozhdéniya.}
\end{quote}

\begin{quote}
My friends wished me a happy birthday.
\end{quote}

Imperfective; perfective \emph{\ru{поздравить}} (\emph{pozdrávit}), \emph{\ru{поздравлю}}. The occasion takes \emph{\ru{с}} + instrumental: \emph{\ru{с} \ru{праздником}!}

\begin{cousinweb}
\emph{\ru{поздравлять}} (EtymDB: \textquotedblleft{}to congratulate\textquotedblright{}) is built on \emph{\ru{по}-} and \emph{\ru{здравый}} \textquotedblleft{}healthy\textquotedblright{}, the Church Slavonic shape (EtymDB: Church Slavonic \emph{sŭdravŭ}) of native \emph{\ru{здоровый}}: to wish someone health.
\end{cousinweb}

\begin{grammarlens}[title={What you've built}]
One more word for telling, explaining and describing at B1.
\end{grammarlens}

\subsection*{Your turn}
\begin{itemize}
\raggedright
  \item \emph{Say it:} \emph{pozdravlyát}
  \item \emph{Say it:} \emph{pozdravlyát}, once more
  \item \emph{Say it:} a sentence of your own with \emph{pozdravlyát}
  \item \emph{Recall:} say the Russian for solemn, then the Russian for festive, then say \emph{pozdravlyát} again
\end{itemize}

\begin{tcolorbox}[breakable,colback=teal!4,colframe=teal!35!black,title={Before you move on}]
What does \ru{поздравлять} mean? (\textquotedblleft{}To congratulate\textquotedblright{}.) Say it once more.
\end{tcolorbox}
\section[\texorpdfstring{\ru{поздравление} (pozdravléniye)}{поздравление (pozdravléniye)}]{\ru{поздравление} (pozdravléniye) — congratulations, a greeting}

\label{lesson:RU-C312-pozdravleniye}



\begin{quote}
Before the new one: say the Russian for festive, then the Russian for to congratulate.
\end{quote}

\subsection*{You'll want to know: \ru{поздравление}}
\textbf{\ru{поздравление}} (\emph{pozdravléniye}) — \textquotedblleft{}congratulations, a greeting\textquotedblright{}.

\begin{quote}
\ru{Она} \ru{получила} \ru{много} \ru{поздравлений}.
\end{quote}

\begin{quote}
\emph{Oná poluchíla mnógo pozdravléniy.}
\end{quote}

\begin{quote}
She received a lot of greetings.
\end{quote}

Neuter.

\begin{cousinweb}
\emph{\ru{поздравление}} is built on \emph{\ru{поздравлять}} \textquotedblleft{}to congratulate\textquotedblright{}, from \emph{\ru{здравый}} \textquotedblleft{}healthy\textquotedblright{}, the Church Slavonic shape of \emph{\ru{здоровый}}.
\end{cousinweb}

\begin{grammarlens}[title={What you've built}]
One more word for telling, explaining and describing at B1.
\end{grammarlens}

\subsection*{Your turn}
\begin{itemize}
\raggedright
  \item \emph{Say it:} \emph{pozdravléniye}
  \item \emph{Say it:} \emph{pozdravléniye}, once more
  \item \emph{Say it:} a sentence of your own with \emph{pozdravléniye}
  \item \emph{Recall:} say the Russian for festive, then the Russian for to congratulate, then say \emph{pozdravléniye} again
\end{itemize}

\begin{tcolorbox}[breakable,colback=teal!4,colframe=teal!35!black,title={Before you move on}]
What does \ru{поздравление} mean? (\textquotedblleft{}Congratulations, a greeting\textquotedblright{}.) Say it once more.
\end{tcolorbox}
\section[\texorpdfstring{\ru{пожелание} (pozhelániye)}{пожелание (pozhelániye)}]{\ru{пожелание} (pozhelániye) — a wish}

\label{lesson:RU-C312-pozhelaniye}



\begin{quote}
Before the new one: say the Russian for to congratulate, then the Russian for congratulations.
\end{quote}

\subsection*{You'll want to know: \ru{пожелание}}
\textbf{\ru{пожелание}} (\emph{pozhelániye}) — \textquotedblleft{}a wish\textquotedblright{}.

\begin{quote}
\ru{Бабушка} \ru{написала} \ru{добрые} \ru{пожелания} \ru{в} \ru{открытке}.
\end{quote}

\begin{quote}
\emph{Bábushka napisála dóbryye pozhelániya v otkrýtke.}
\end{quote}

\begin{quote}
Grandma wrote kind wishes in the card.
\end{quote}

Neuter. \emph{\ru{С} \ru{наилучшими} \ru{пожеланиями}} — with best wishes.

\begin{cousinweb}
\emph{\ru{пожелание}} is built on \emph{\ru{пожелать}}, from \emph{\ru{желать}} \textquotedblleft{}to wish\textquotedblright{}, inherited from Proto-Slavic \emph{želati} (EtymDB).
\end{cousinweb}

\begin{grammarlens}[title={What you've built}]
One more word for telling, explaining and describing at B1.
\end{grammarlens}

\subsection*{Your turn}
\begin{itemize}
\raggedright
  \item \emph{Say it:} \emph{pozhelániye}
  \item \emph{Say it:} \emph{pozhelániye}, once more
  \item \emph{Say it:} a sentence of your own with \emph{pozhelániye}
  \item \emph{Recall:} say the Russian for to congratulate, then the Russian for congratulations, then say \emph{pozhelániye} again
\end{itemize}

\begin{tcolorbox}[breakable,colback=teal!4,colframe=teal!35!black,title={Before you move on}]
What does \ru{пожелание} mean? (\textquotedblleft{}A wish\textquotedblright{}.) Say it once more.
\end{tcolorbox}
\section[\texorpdfstring{\ru{желать} (zhelát)}{желать (zhelát)}]{\ru{желать} (zhelát) — to wish}

\label{lesson:RU-C312-zhelat}



\begin{quote}
Before the new one: say the Russian for congratulations, then the Russian for a wish.
\end{quote}

\subsection*{You'll want to know: \ru{желать}}
\textbf{\ru{желать}} (\emph{zhelát}) — \textquotedblleft{}to wish\textquotedblright{}.

\begin{quote}
\ru{Желаю} \ru{тебе} \ru{счастья}!
\end{quote}

\begin{quote}
\emph{Zheláyu tebé schástya!}
\end{quote}

\begin{quote}
I wish you happiness!
\end{quote}

Imperfective; perfective \emph{\ru{пожелать}} (\emph{pozhelát}). The person is in the dative, the wish in the genitive: \emph{\ru{желаю} \ru{удачи}}.

\begin{cousinweb}
\emph{\ru{желать}} is inherited from Proto-Slavic \emph{želati} \textquotedblleft{}to wish, to long for\textquotedblright{} (EtymDB).
\end{cousinweb}

\begin{grammarlens}[title={What you've built}]
One more word for telling, explaining and describing at B1.
\end{grammarlens}

\subsection*{Your turn}
\begin{itemize}
\raggedright
  \item \emph{Say it:} \emph{zhelát}
  \item \emph{Say it:} \emph{zhelát}, once more
  \item \emph{Say it:} a sentence of your own with \emph{zhelát}
  \item \emph{Recall:} say the Russian for congratulations, then the Russian for a wish, then say \emph{zhelát} again
\end{itemize}

\begin{tcolorbox}[breakable,colback=teal!4,colframe=teal!35!black,title={Before you move on}]
What does \ru{желать} mean? (\textquotedblleft{}To wish\textquotedblright{}.) Say it once more.
\end{tcolorbox}
\section[\texorpdfstring{\ru{тост} (tost)}{тост (tost)}]{\ru{тост} (tost) — a toast (with drinks)}

\label{lesson:RU-C312-tost}



\begin{quote}
Before the new one: say the Russian for a wish, then the Russian for to wish.
\end{quote}

\subsection*{You'll want to know: \ru{тост}}
\textbf{\ru{тост}} (\emph{tost}) — \textquotedblleft{}a toast (with drinks)\textquotedblright{}.

\begin{quote}
\ru{Отец} \ru{произнёс} \ru{тост} \ru{за} \ru{молодых}.
\end{quote}

\begin{quote}
\emph{Otéts proiznyós tost za molodýkh.}
\end{quote}

\begin{quote}
Father proposed a toast to the newlyweds.
\end{quote}

\emph{\ru{Поднять} \ru{тост} \ru{за}} + accusative. Also toasted bread.

\begin{cousinweb}
\emph{\ru{тост}} was borrowed from English \emph{toast} (EtymDB), from Old French \emph{toster} \textquotedblleft{}to roast\textquotedblright{}, from Latin \emph{tostus} \textquotedblleft{}roasted\textquotedblright{}, the participle of \emph{torrēre} \textquotedblleft{}to parch\textquotedblright{}.
\end{cousinweb}

\begin{grammarlens}[title={What you've built}]
One more word for telling, explaining and describing at B1.
\end{grammarlens}

\subsection*{Your turn}
\begin{itemize}
\raggedright
  \item \emph{Say it:} \emph{tost}
  \item \emph{Say it:} \emph{tost}, once more
  \item \emph{Say it:} a sentence of your own with \emph{tost}
  \item \emph{Recall:} say the Russian for a wish, then the Russian for to wish, then say \emph{tost} again
\end{itemize}

\begin{tcolorbox}[breakable,colback=teal!4,colframe=teal!35!black,title={Before you move on}]
What does \ru{тост} mean? (\textquotedblleft{}A toast (with drinks)\textquotedblright{}.) Say it once more.
\end{tcolorbox}
